% Margin sparkline: escalation threshold u*(delta) = delta / (1 + delta), w=1.
% Domain delta in [0, 0.50]; clears cost above threshold.
% Fits inside the 1.3in book margin column.
% Reader question: how much of the validation signal u must an escalation
%   clear before it is worth its cost, as the debit delta grows?
% Claim: the threshold rises with delta but stays well under the 0.333
%   ceiling across the whole worked range; the chapter's own hand-worked
%   point is delta=0.05, where u* = 0.0476.
% Source: sec:escalation, Eq. u* = delta*w / (1 + delta*w); the worked
%   calculation immediately below this sparkline's \marginnote call.
\begingroup\makeatletter
\ifdefined\pd@typesetpicturebox
  \let\pgfsys@typesetpicturebox\pd@typesetpicturebox
\fi
\begin{tikzpicture}[pd figure,x=115pt,y=200pt]
  % Baseline and upper reference: neutral, never a drawn edge.
  \draw[pd hairline] (0,0) -- (0.50,0);
  \draw[pd guide] (0,0.333) -- (0.50,0.333);
  % Escalation threshold curve u*(delta) = delta / (1 + delta).
  \draw[pd legible,line width=.9pt] plot[domain=0:0.50,samples=31]
    (\x,{\x / (1 + \x)});
  % The chapter's own hand-worked point (delta=0.05) is the focus; the
  % second point (delta=0.30) is a plain reference further up the curve.
  \node[pd focus datum] at (0.05,{0.05/1.05}) {};
  \node[pd legible datum] at (0.30,{0.30/1.30}) {};
  % Labels: end values, the formula, and the worked range.
  \node[pd label,anchor=east] at (0,0) {$0$};
  \node[pd label,anchor=south west] at (0.50,0.333) {$.33$};
  \node[pd direct label,anchor=south east] at (0.49,0.012) {$\delta/(1+\delta)$};
  \node[pd label,anchor=north] at (0.25,-0.035) {debit $\delta: 0\longrightarrow .50$};
\end{tikzpicture}
\makeatother\endgroup
